(1) 非相互作用二维Dirac粒子具有对称的导带和价带 \(E_\pm(\boldsymbol k)=\pm\hbar v|\boldsymbol k|\)，速度常数 \(v>0\)、内部简并度 \(g\)。以中性化满价带为能量零参考，温度 \(T>0\)、化学势 \(\mu_{\rm ch}=0\) 时，电子与空穴热激发彼此对称。求每单位面积正能量态密度 \(D(E)\)、热电子数密度 \(n_e\)、热空穴 \(n_h\)。
(2) 通过显式积分 \( \int_0^\infty x^2/(e^x+1)\,dx=3\zeta(3)/2\)，求相对于填满价带的总内能密度 \(u(T)\) 和定面积热容 \(c_A=(\partial u/\partial T)_A\)，核对低温标度与量纲。
(3) 由线性色散的尺度变换或巨配分函数，证明热激发压力 \(P=u/2\)，求熵面密度 \(s\)。解释这不是三维理想气体 \(P=2u/3\)，并检验 \(T\to0\) 的电中性基态极限。
(4) 允许化学势施加小偏移 \(\mu_{\rm ch}\)，写净电荷相关粒子数 \(n_{\rm net}=n_e-n_h\) 的积分表达式。求零化学势的静态压缩响应 \(\chi=\partial n_{\rm net}/\partial\mu_{\rm ch}|_0\) 和量子电容面密度 \(C_Q=e^2\chi\)，说明该非零热量子电容与零温中性点的区别及忽略相互作用的边界。
